\section{Conclusion}
\label{sec:conclusion}

This report followed one geometric path from an author's implicit shape to
a shaded pixel.  The CPU evaluates the SDF, builds an adaptive octree of
corner samples, simplifies it into a LoD ladder, and flattens a chunk's
neighbourhood into compact tables.  The GPU walks the twelve cube edges,
stitches LoD boundaries, solves one dual vertex per surface cell, and
writes the final 64-byte vertices and chunk-local indices directly into
the renderer's persistent pools.  The host sees counters, slots, and
versions—not vertices.

The key ideas are:
\begin{itemize}
\item \textbf{One Cartesian producer.}  The geometry pass is the only runtime code
  that writes a vertex.  This eliminates an entire class of
  CPU/GPU divergence and makes the offline exporter a meaningful
  reference.
\item \textbf{Root-consistent flattening.}  The GPU table is a
  serialisable snapshot of the exact neighbourhood the CPU walker would
  traverse, including the one-cell margin for coarse neighbours.
\item \textbf{Break-based stitching.}  The edge is split at every cell
  boundary, not at a fixed table of configurations.  Two chunks resolve
  the same edge to the same owner and the same sub-segments.
\item \textbf{Span publishing.}  A chunk's geometry is a reserved range in
  a large pool, published only after the geometry fence and freed only
  after the cull frames that could reference it have retired.
\item \textbf{One hierarchy for everything.}  Raster, shadows, ray
  tracing, vegetation, and fire all consume the same published span or the
  same node metadata.
\end{itemize}

The result is not a toy mesher.  It streams an adaptive, crack-free
surface from a mutable SDF octree into a real-time renderer without a
vertex readback, and it keeps the host's work proportional to the number
of chunks and nodes, not to the number of triangles.  The measured costs
and the invariants in this report are the contract a future change must
preserve.
